\chapter{Folate}

\section*{Definition and Overview}
Folate, also called vitamin B9 or folic acid in its synthetic form, is a water-soluble B vitamin essential for the production of red blood cells, the synthesis of DNA, and normal growth and development. The body needs folate throughout life, but it is especially important during pregnancy, when it prevents serious birth defects of the brain and spine (neural tube defects). Folate is found naturally in green leafy vegetables, legumes, and some fruits, and folic acid is added to many fortified foods and supplements. Folate deficiency causes anaemia and can harm a developing baby, but it is easily prevented and treated. This chapter explains what folate does, who is at risk of deficiency, and how to ensure adequate intake.

\section*{Epidemiology}
Folate deficiency is relatively uncommon in Australia since the introduction of mandatory folic acid fortification of bread flour, which reduced neural tube defects by around 14\% and raised population folate levels substantially. However, certain groups remain at risk: women of childbearing age who do not take supplements, people with poor diets, people with malabsorption conditions such as coeliac disease, people taking medicines that interfere with folate, heavy drinkers, and some people with anaemia. Neural tube defects affect around 1--2 per 1,000 pregnancies in Australia without supplementation, making preconception folate a priority in public health.

\section*{Aetiology and Pathophysiology}
\begin{itemize}
    \item \textbf{Role of folate:} a cofactor in the production of DNA and red blood cells, and in amino acid metabolism
    \item \textbf{Deficiency causes:} inadequate dietary intake, malabsorption (coeliac disease, inflammatory bowel disease), increased demand (pregnancy, growth, haemolysis), and medicines including methotrexate, sulfasalazine, and some anticonvulsants
    \item \textbf{Alcohol:} interferes with folate absorption and increases its loss
    \item \textbf{Neural tube defects:} in early pregnancy, folate deficiency impairs closure of the developing spinal cord and brain, causing spina bifida and anencephaly
    \item \textbf{Megaloblastic anaemia:} without folate, red cell production fails, and the large, immature red cells characteristic of the condition appear in the blood
    \item \textbf{Homocysteine:} folate deficiency raises homocysteine, an amino acid linked to cardiovascular risk
    \item \textbf{Folic acid versus folate:} folic acid is the synthetic form used in supplements and fortified foods; it is more stable and better absorbed than natural folate
\end{itemize}

\section*{Clinical Features}
\begin{itemize}
    \item In early deficiency, there may be no symptoms
    \item Anaemia: fatigue, weakness, pallor, breathlessness on exertion, and palpitations
    \item Sore, red, smooth tongue (glossitis)
    \item Mouth ulcers and angular cheilitis at the corners of the mouth
    \item Irritability, poor concentration, and depressed mood
    \item Reduced appetite, weight loss, and diarrhoea
    \item In pregnancy: neural tube defects in the baby when intake is inadequate in the first weeks
    \item Neurological symptoms are uncommon with folate deficiency (unlike B12 deficiency), but can occur
\end{itemize}

\section*{Differential Diagnosis}
Anaemia and its symptoms have several causes that must be distinguished.
\begin{itemize}
    \item Vitamin B12 deficiency, which causes a similar megaloblastic anaemia plus neurological symptoms
    \item Iron deficiency anaemia
    \item Anaemia of chronic disease
    \item Hypothyroidism, which can also cause fatigue
    \item Coeliac disease, which causes malabsorption of folate and other nutrients
    \item The key test distinction: B12 deficiency must be excluded before treating with folate alone, because folate can mask B12 deficiency and allow its nerve damage to progress
\end{itemize}

\section*{When to Seek Medical Care}
Anyone with persistent fatigue, a sore tongue, or mouth ulcers should see a doctor for blood tests. Women planning pregnancy or in early pregnancy need adequate folate, and people on medicines that interfere with folate should discuss supplementation with their doctor.

\textbf{Call 000 (or local emergency services) for:}
\begin{itemize}
    \item Severe breathlessness, chest pain, or fainting from profound anaemia
    \item Sudden weakness or neurological symptoms
    \item Any emergency symptoms
\end{itemize}

\section*{Diagnosis and Investigations}
\begin{itemize}
    \item \textbf{Full blood count:} shows macrocytic (large red cell) anaemia in deficiency
    \item \textbf{Serum folate:} low in deficiency, though levels can fluctuate with recent intake
    \item \textbf{Red cell folate:} reflects longer-term folate status
    \item \textbf{Vitamin B12 measurement:} essential to exclude B12 deficiency before treatment
    \item \textbf{Investigation of the cause:} dietary review, and testing for coeliac disease or other malabsorption where suspected
    \item \textbf{Homocysteine:} raised in folate deficiency, sometimes measured
\end{itemize}

\section*{Management}
\begin{itemize}
    \item \textbf{Dietary folate:} green leafy vegetables, legumes, citrus fruits, eggs, and fortified breads and cereals
    \item \textbf{Folic acid supplements:} 500 micrograms daily is the standard treatment for deficiency, and larger doses for specific conditions
    \item \textbf{Before and during pregnancy:} 500 micrograms daily from at least one month before conception until the end of the first trimester, with higher doses (5 mg) for women at increased risk
    \item \textbf{Treat the cause:} gluten-free diet for coeliac disease, and medication review where medicines interfere
    \item \textbf{Treatment of anaemia:} blood counts are repeated to confirm recovery
    \item \textbf{Ongoing supplements:} for people with continuing malabsorption or long-term medication use
    \item \textbf{Alcohol reduction:} for heavy drinkers with deficiency
\end{itemize}

\section*{Complications}
\begin{itemize}
    \item Neural tube defects (spina bifida, anencephaly) in babies of folate-deficient mothers
    \item Megaloblastic anaemia with fatigue and reduced function
    \item Progression of undiagnosed B12 deficiency if folate is given alone
    \item Adverse outcomes in pregnancy, including low birth weight
    \item In severe, chronic deficiency, cognitive and mood effects
\end{itemize}

\section*{Prognosis and Outcomes}
Folate deficiency responds rapidly to treatment: symptoms improve within days to weeks, and the blood count normalises over one to two months. The prevention of neural tube defects is highly effective when supplementation starts before conception, which is why it is recommended for all women planning pregnancy. With a balanced diet and targeted supplementation for at-risk groups, deficiency is almost entirely preventable, and its complications are avoided. The long-term outlook for treated folate deficiency is excellent.

\section*{Prevention and Self-Care}
\begin{itemize}
    \item Eat folate-rich foods: spinach, broccoli, legumes, citrus, and fortified cereals
    \item Take 500 micrograms of folic acid daily from one month before conception until 12 weeks of pregnancy
    \item Women at higher risk, including those with a previous neural tube defect or on certain medicines, should take 5 mg daily on medical advice
    \item Limit alcohol, which interferes with folate
    \item Ask a doctor about folate supplementation if taking methotrexate or other interfering medicines
    \item Do not start folate without checking B12 when treating unexplained anaemia
    \item Attend routine health checks, including blood counts where indicated
\end{itemize}

\section*{Sources and Further Reading}
The following reputable sources provide further information for healthcare students and patients seeking reliable, current advice about folate.
\begin{itemize}
    \item Healthdirect Australia. \textit{Folate and folic acid.} https://www.healthdirect.gov.au/folate
    \item Food Standards Australia New Zealand. \textit{Folic acid fortification.} https://www.foodstandards.gov.au/
    \item Nutrition Australia. \textit{Folate.} https://nutritionaustralia.org/
    \item NHS. \textit{Vitamin B and folic acid.} https://www.nhs.uk/conditions/vitamins-and-minerals/
    \item Mayo Clinic. \textit{Folate deficiency.} https://www.mayoclinic.org/
    \item MedlinePlus. \textit{Folate.} https://medlineplus.gov/
\end{itemize}
